 

%请使用xelatex编译
\documentclass[11pt,addpoints,answers]{exam}


% required macros -- get header.tex file from course web page
% !Mode:: "TeX:UTF-8"
% !TEX root=./hw1.tex
\usepackage{xeCJK}    % 写中文要用到
\usepackage{tikz}
\usepackage{amsmath,amsfonts,amssymb,amsthm}
\usepackage{fullpage}
\usepackage{times}
\usepackage{hyperref}
\usepackage{pdfsync}
\usepackage{microtype}
\usepackage{color}
\usepackage{cleveref}
\crefformat{footnote}{#2\footnotemark[#1]#3}
\definecolor{light-gray}{gray}{0.5}

\renewcommand{\tablename}{表}
\renewcommand{\figurename}{图}
\renewcommand{\proof}{证明}

% \theoremstyle{plain}
% \newtheorem{theorem}{定理~}
% \newtheorem{lemma}{引理~}
% \newtheorem{axiom}{公理~}
% \newtheorem{proposition}{命题~}
% \newtheorem{prop}{性质~}
% \newtheorem{corollary}{推论~}
% \newtheorem{conclusion}{结论~}
% \newtheorem{definition}{定义~}
% \newtheorem{conjecture}{猜想~}
% \newtheorem{example}{例~}
% \newtheorem{remark}{注~}
% \newtheorem{algorithm}{算法~}


%%% BLACKBOARD SYMBOLS

% \newcommand{\C}{\ensuremath{\mathbb{C}}}
\newcommand{\D}{\ensuremath{\mathbb{D}}}
\newcommand{\F}{\ensuremath{\mathbb{F}}}
% \newcommand{\G}{\ensuremath{\mathbb{G}}}
\newcommand{\J}{\ensuremath{\mathbb{J}}}
\newcommand{\N}{\ensuremath{\mathbb{N}}}
\newcommand{\Q}{\ensuremath{\mathbb{Q}}}
\newcommand{\R}{\ensuremath{\mathbb{R}}}
\newcommand{\T}{\ensuremath{\mathbb{T}}}
\newcommand{\Z}{\ensuremath{\mathbb{Z}}}
\newcommand{\QR}{\ensuremath{\mathbb{QR}}}

\newcommand{\Zt}{\ensuremath{\Z_t}}
\newcommand{\Zp}{\ensuremath{\Z_p}}
\newcommand{\Zq}{\ensuremath{\Z_q}}
\newcommand{\ZN}{\ensuremath{\Z_N}}
\newcommand{\Zps}{\ensuremath{\Z_p^*}}
\newcommand{\ZNs}{\ensuremath{\Z_N^*}}
\newcommand{\JN}{\ensuremath{\J_N}}
\newcommand{\QRN}{\ensuremath{\QR_{N}}}
\newcommand{\QRp}{\ensuremath{\QR_{p}}}

%%% THEOREM COMMANDS

\theoremstyle{plain}            % following are "theorem" style
\newtheorem{theorem}{定理}[section]
\newtheorem{lemma}[theorem]{引理}
\newtheorem{corollary}[theorem]{推论}
\newtheorem{proposition}[theorem]{命题}
\newtheorem{claim}[theorem]{声明}
\newtheorem{fact}[theorem]{事实}

\theoremstyle{definition}       % following are def style
\newtheorem{definition}[theorem]{定义}
\newtheorem{conjecture}[theorem]{推测}
\newtheorem{example}[theorem]{例}
\newtheorem{protocol}[theorem]{协议}

\theoremstyle{remark}           % following are remark style
\newtheorem{remark}[theorem]{注}
\newtheorem{note}[theorem]{注}
\newtheorem{exercise}[theorem]{练习}

% equation numbering style
\numberwithin{equation}{section}

%%% GENERAL COMPUTING

\newcommand{\bit}{\ensuremath{\set{0,1}}}
\newcommand{\pmone}{\ensuremath{\set{-1,1}}}

% asymptotics
\DeclareMathOperator{\poly}{poly}
\DeclareMathOperator{\polylog}{polylog}
\DeclareMathOperator{\negl}{negl}
\newcommand{\Otil}{\ensuremath{\tilde{O}}}

% probability/distribution stuff
\DeclareMathOperator*{\E}{E}
\DeclareMathOperator*{\Var}{Var}

% sets in calligraphic type
\newcommand{\calA}{\ensuremath{\mathcal{A}}}
\newcommand{\calD}{\ensuremath{\mathcal{D}}}
\newcommand{\calF}{\ensuremath{\mathcal{F}}}
\newcommand{\calH}{\ensuremath{\mathcal{H}}}
\newcommand{\calK}{\ensuremath{\mathcal{K}}}
\newcommand{\calM}{\ensuremath{\mathcal{M}}}
\newcommand{\calX}{\ensuremath{\mathcal{X}}}
\newcommand{\calY}{\ensuremath{\mathcal{Y}}}

% types of indistinguishability
\newcommand{\compind}{\ensuremath{\stackrel{c}{\approx}}}
\newcommand{\statind}{\ensuremath{\stackrel{s}{\approx}}}
\newcommand{\perfind}{\ensuremath{\equiv}}

% font for general-purpose algorithms
\newcommand{\algo}[1]{\ensuremath{\mathsf{#1}}}
% font for general-purpose computational problems
\newcommand{\problem}[1]{\ensuremath{\mathsf{#1}}}
% font for complexity classes
\newcommand{\class}[1]{\ensuremath{\mathsf{#1}}}

% complexity classes and languages
\renewcommand{\P}{\class{P}}
\newcommand{\BPP}{\class{BPP}}
\newcommand{\NP}{\class{NP}}
\newcommand{\coNP}{\class{coNP}}
\newcommand{\AM}{\class{AM}}
\newcommand{\coAM}{\class{coAM}}
\newcommand{\IP}{\class{IP}}

%%% "LEFT-RIGHT" PAIRS OF SYMBOLS

% inner product
\newcommand{\inner}[1]{\langle{#1}\rangle}
\newcommand{\innerfit}[1]{\left\langle{#1}\right\rangle}
% absolute value
\newcommand{\abs}[1]{\lvert{#1}\rvert}
\newcommand{\absfit}[1]{\left\lvert{#1}\right\rvert}
% a set
\newcommand{\set}[1]{\{{#1}\}}
\newcommand{\setfit}[1]{\left\{{#1}\right\}}
% parens
\newcommand{\parens}[1]{({#1})}
\newcommand{\parensfit}[1]{\left({#1}\right)}
% tuple = alias for parens
\newcommand{\tuple}[1]{\parens{#1}}
\newcommand{\tuplefit}[1]{\parensfit{#1}}
% square brackets
\newcommand{\bracks}[1]{[{#1}]}
\newcommand{\bracksfit}[1]{\left[{#1}\right]}
% rounding off
\newcommand{\round}[1]{\lfloor{#1}\rceil}
% floor function
\newcommand{\floor}[1]{\lfloor{#1}\rfloor}
% ceiling function
\newcommand{\ceil}[1]{\lceil{#1}\rceil}
% length of a string
\newcommand{\len}[1]{\lvert{#1}\rvert}
\newcommand{\lenfit}[1]{\left\lvert{#1}\right\rvert}
% length of some vector, element
\newcommand{\length}[1]{\lVert{#1}\rVert}
\newcommand{\lengthfit}[1]{\left\lVert{#1}\right\rVert}

%%% CRYPTO-RELATED NOTATION

% KEYS AND RELATED

\newcommand{\key}[1]{\ensuremath{#1}}

\newcommand{\pk}{\key{pk}}
\newcommand{\vk}{\key{vk}}
\newcommand{\sk}{\key{sk}}
\newcommand{\mpk}{\key{mpk}}
\newcommand{\msk}{\key{msk}}
\newcommand{\fk}{\key{fk}}
\newcommand{\id}{id}
\newcommand{\keyspace}{\ensuremath{\mathcal{K}}}
\newcommand{\msgspace}{\ensuremath{\mathcal{M}}}
\newcommand{\ctspace}{\ensuremath{\mathcal{C}}}
\newcommand{\tagspace}{\ensuremath{\mathcal{T}}}
\newcommand{\idspace}{\ensuremath{\mathcal{ID}}}

\newcommand{\concat}{\ensuremath{\|}}

% GAMES

% advantage
\newcommand{\advan}{\ensuremath{\mathbf{Adv}}}

% different attack models
\newcommand{\attack}[1]{\ensuremath{\text{#1}}}

\newcommand{\atk}{\attack{atk}} % dummy attack
\newcommand{\indcpa}{\attack{ind-cpa}}
\newcommand{\indcca}{\attack{ind-cca}}
\newcommand{\anocpa}{\attack{ano-cpa}} % anonymous
\newcommand{\anocca}{\attack{ano-cca}}
\newcommand{\euacma}{\attack{eu-acma}} % forgery: adaptive chosen-message
\newcommand{\euscma}{\attack{eu-scma}} % forgery: static chosen-message
\newcommand{\suacma}{\attack{su-acma}} % strongly unforgeable

% ADVERSARIES
\newcommand{\attacker}[1]{\ensuremath{\mathcal{#1}}}

\newcommand{\Adv}{\attacker{A}}
\newcommand{\AdvA}{\attacker{A}}
\newcommand{\AdvB}{\attacker{B}}
\newcommand{\Dist}{\attacker{D}}
\newcommand{\Sim}{\attacker{S}}
\newcommand{\Ora}{\attacker{O}}
\newcommand{\Inv}{\attacker{I}}
\newcommand{\For}{\attacker{F}}

% CRYPTO SCHEMES

\newcommand{\scheme}[1]{\ensuremath{\text{#1}}}

% pseudorandom stuff
\newcommand{\prg}{\algo{PRG}}
\newcommand{\prf}{\algo{PRF}}
\newcommand{\prp}{\algo{PRP}}

% symmetric-key cryptosystem
\newcommand{\skc}{\scheme{SKC}}
\newcommand{\skcgen}{\algo{Gen}}
\newcommand{\skcenc}{\algo{Enc}}
\newcommand{\skcdec}{\algo{Dec}}

% public-key cryptosystem
\newcommand{\pkc}{\scheme{PKC}}
\newcommand{\pkcgen}{\algo{Gen}}
\newcommand{\pkcenc}{\algo{Enc}} % can also use \kemenc and \kemdec
\newcommand{\pkcdec}{\algo{Dec}}

% digital signatures
\newcommand{\sig}{\scheme{SIG}}
\newcommand{\siggen}{\algo{Gen}}
\newcommand{\sigsign}{\algo{Sign}}
\newcommand{\sigver}{\algo{Ver}}

% message authentication code
\newcommand{\mac}{\scheme{MAC}}
\newcommand{\macgen}{\algo{Gen}}
\newcommand{\mactag}{\algo{Tag}}
\newcommand{\macver}{\algo{Ver}}

% key-encapsulation mechanism
\newcommand{\kem}{\scheme{KEM}}
\newcommand{\kemgen}{\algo{Gen}}
\newcommand{\kemenc}{\algo{Encaps}}
\newcommand{\kemdec}{\algo{Decaps}}

% identity-based encryption
\newcommand{\ibe}{\scheme{IBE}}
\newcommand{\ibesetup}{\algo{Setup}}
\newcommand{\ibeext}{\algo{Ext}}
\newcommand{\ibeenc}{\algo{Enc}}
\newcommand{\ibedec}{\algo{Dec}}

% hierarchical IBE (as key encapsulation)
\newcommand{\hibe}{\scheme{HIBE}}
\newcommand{\hibesetup}{\algo{Setup}}
\newcommand{\hibeext}{\algo{Extract}}
\newcommand{\hibeenc}{\algo{Encaps}}
\newcommand{\hibedec}{\algo{Decaps}}

% binary tree encryption (as key encapsulation)
\newcommand{\bte}{\scheme{BTE}}
\newcommand{\btesetup}{\algo{Setup}}
\newcommand{\bteext}{\algo{Extract}}
\newcommand{\bteenc}{\algo{Encaps}}
\newcommand{\btedec}{\algo{Decaps}}

% trapdoor functions
\newcommand{\tdf}{\scheme{TDF}}
\newcommand{\tdfgen}{\algo{Gen}}
\newcommand{\tdfeval}{\algo{Eval}}
\newcommand{\tdfinv}{\algo{Invert}}
\newcommand{\tdfver}{\algo{Ver}}

%%% PROTOCOLS

\newcommand{\out}{\text{out}}
\newcommand{\view}{\text{view}}

%%% COMMANDS FOR LECTURES/HOMEWORKS

\newcommand{\lecheader}{%
  \chead{\large \textbf{Lecture \lecturenum\\\lecturetopic}}

  \lhead{\small
    \textbf{\href{http://bhpan.buaa.edu.cn}{研究生课程：计算复杂性}\\2023年秋季}}

  \rhead{\small \textbf{主讲: 张宗洋
      }

  \setlength{\headheight}{20pt}
  \setlength{\headsep}{16pt}
}}

\newcommand{\hwheader}{%
  \chead{\Large \textbf{作业 \hwnum}}

  \lhead{\small
    \textbf{{计算复杂性}\\2025年秋季}}

  \rhead{\small \textbf{主讲: 张宗洋
      \\姓名: \studentname}}

  \setlength{\headheight}{20pt}
  \setlength{\headsep}{16pt}
  
  \headrule
}

\newcommand{\examheader}{%
  \chead{\Large \textbf{测试 \\ \duedate}}
  
\lhead{\small
    \textbf{\href{http://bhpan.buaa.edu.cn}{计算复杂性}\\2024年秋季}}

  \rhead{\small \textbf{主讲: 张宗洋
      \\姓名: \studentname}}

  \setlength{\headheight}{20pt}
  \setlength{\headsep}{16pt}
  
  \headrule
}
  
	
\newcommand{\hint}[1]{\href{http://www.cims.nyu.edu/~regev/cgi-bin/hints/index.html?#1}{{\textcolor{light-gray}{I need a hint! (ID #1)}}}}

\newcommand{\hinttext}[2]{\href{http://www.cims.nyu.edu/~regev/cgi-bin/hints/index.html?#1}{{\textcolor{light-gray}{#2 (ID #1)}}}}

\newcommand{\hintcost}[2]{\href{http://www.cims.nyu.edu/~regev/cgi-bin/hints/index.html?#1}{{\textcolor{light-gray}{I need a hint for #2 points! (ID #1)}}}}
	
\newcommand{\moreinfo}[1]{\href{http://www.cims.nyu.edu/~regev/cgi-bin/hints/index.html?#1}{{\textcolor{light-gray}{I'm done solving and want to know more! (ID #1)}}}}

% VARIABLES

\newcommand{\hwnum}{4}
\newcommand{\duedate}{12月7日23时59分前} % changing this does not change the
                                % actual due date :)
\newcommand{\studentname}{}

% END OF SUPPLIED VARIABLES

\hwheader                       % execute homework commands

\begin{document}

\pagestyle{head}                % put header on every page
\begin{itemize}
	\item 请将PDF文件发送至邮箱：complexity\_buaa@163.com
	\item 文件命名规则是``学号-姓名-第四次作业.pdf"
	\item 务必独立自主解题
	\item 作业截止时间为 \textbf{\duedate}.
\end{itemize}



\begin{questions}
	\question P132页, 4.3题.
	设$ALL_{\textit{DFA}}=\left\{ \langle A \rangle| A\text{是一个} \textit{DFA} \text{,且}L(A)=\Sigma ^{*} \right\}$。证明$ALL_{\textit{DFA}}$是可判定的。
	\begin{solution}
    
		构造一个图灵机 \(M\)，输入为 DFA \(A\)：
        \begin{enumerate}
            \item 构造 DFA \(B\)，使得 \(L(B) = \overline{L(A)}\)（对 DFA 取补是有效的）。
            \item 检查 \(L(B)\) 是否为空：运行判定器 \(E_{DFA}\)（已知 \(E_{DFA}\) 可判定）。
            \item 若 \(E_{DFA}(B)\) 接受，则说明 \(L(A) = \Sigma^*\)，输出接受；否则输出拒绝。
        \end{enumerate}
        由于每一步都可计算，因此 \(ALL_{DFA}\) 可判定。
	\end{solution}
	
	
	\question P133页, 4.4题.
	设$A\epsilon_{\textit{CFG}}=\left\{ \langle G \rangle| G\text{是一个生成} \epsilon \text{的} \textit{CFG} \right\}$。证明$A\epsilon_{ \textit{CFG}}$是可判定的。
	\begin{solution}
    
    对 CFG \(G\)，构造如下判定过程：
    \begin{enumerate}
        \item 初始化集合 $\text{Nullable} = \varnothing$。
        \item 重复以下步骤直到没有新符号加入 $\text{Nullable}$：
        \item 对每条产生式 $A \to \alpha \in R$：
        \begin{itemize}
            \item 若 $\alpha = \epsilon$，则将 $A$ 加入 $\text{Nullable}$。
            \item 若 $\alpha = B_1 B_2 \dots B_k$（$k \ge 1$，每个 $B_i \in V \cup \Sigma$），并且每个 $B_i$ 都在 $\text{Nullable}$ 中（注意：终结符不在 $\text{Nullable}$ 中），则将 $A$ 加入 $\text{Nullable}$。
        \end{itemize}
        \item 如果 $S \in \text{Nullable}$，则接受；否则拒绝。
    \end{enumerate}
    该算法在多项式时间内完成，因此 \(A\epsilon_{CFG}\) 可判定。
	\end{solution}  
	
	\question P133页, 4.5题.
	设$E_{\textit{TM}}=\left\{ \langle M \rangle| M\text{是一个图灵机} \text{且} L(M)= \varnothing \right\}$。证明$E_{\textit{TM}}$的补$\overline{E_{\textit{TM}}}$是图灵可识别的。
	\begin{solution}
    
		构造图灵机 $R$ 如下：
        $R =$ “对于输入 $\langle M \rangle$，其中 $M$ 是一个图灵机：
        
        对于 $i = 1,2,3,\dots$ 重复以下步骤：
        \begin{enumerate}
            \item 对于所有长度不超过 $i$ 的字符串 $s$（按某种固定顺序枚举）：
            \item 在 $M$ 上模拟运行 $i$ 步，输入为 $s$。
            \item 如果在模拟过程中 $M$ 接受了 $s$，则接受 $\langle M \rangle$。”
        \end{enumerate}
	\end{solution}
	
	\question P133页, 4.14题.
	设$C=\{\langle G,x\rangle|G\text{是一个CFG，}x\text{是某个}y\in L(G)\text{的子串}\}$，证明$C$是可判定的。（提示：使用$E_\text{CFG}$的判定器是个好方法）
	\begin{solution}
		考虑如下判定器 \(D\) 的构造：

        \(D =\) “对于输入 \(\langle G, x \rangle\)，其中 \(G = (V, \Sigma, R, S)\) 是一个 CFG，\(x \in \Sigma^*\)：
        \begin{enumerate}
            \item 构造 CFG \(G'\) 使得 \(L(G') = \Sigma^* x \Sigma^*\)（即所有包含 \(x\) 作为子串的字符串组成的正则语言）。  
                  具体构造：设 \(x = a_1 a_2 \dots a_n\)，可以构造右线性文法：
                  \begin{align*}
                      S' &\to A_1 x A_2, \\
                      A_1 &\to \epsilon \mid \Sigma A_1, \\
                      A_2 &\to \epsilon \mid \Sigma A_2,
                  \end{align*}
                  或者直接用正则表达式转 CFG 的标准方法。
            \item 构造 CFG \(G''\) 使得 \(L(G'') = L(G) \cap L(G')\)。  
                  因为 \(L(G')\) 是正则语言，而 CFG 与正则语言的交仍是 CFG，且该构造可以在算法上实现（例如，将 \(G\) 转化为等价 PDA，与 \(L(G')\) 的 DFA 做乘积，再转回 CFG；或者直接使用已知的闭包性质构造算法）。
            \item 在 \(G''\) 上运行 \(E_{\text{CFG}}\) 的判定器（检查 \(L(G'')\) 是否为空）。
            \item 如果 \(E_{\text{CFG}}\) 判定器输出“非空”，则接受 \(\langle G,x\rangle\)；否则拒绝。”
        \end{enumerate}
                
        \textbf{正确性分析}：
        \begin{itemize}
            \item 若 \(\langle G, x \rangle \in C\)，则存在字符串 \(w \in L(G)\) 且 \(x\) 是 \(w\) 的子串，即 \(w \in L(G')\)。于是 \(w \in L(G) \cap L(G') = L(G'')\)，故 \(L(G'') \neq \varnothing\)，\(E_{\text{CFG}}\) 判定器输出“非空”，因此 \(D\) 接受。
            \item 若 \(\langle G, x \rangle \notin C\)，则任何 \(w \in L(G)\) 都不包含 \(x\) 作为子串，因此 \(L(G) \cap L(G') = \varnothing\)，即 \(L(G'') = \varnothing\)，\(E_{\text{CFG}}\) 判定器输出“空”，因此 \(D\) 拒绝。
        \end{itemize}
        
        因为 \(E_{\text{CFG}}\) 是可判定的，并且构造 \(G'\) 与 \(G''\) 的过程都是有限步骤可完成的，所以 \(D\) 是一个判定器。  
        故 \(C\) 是可判定的。
	\end{solution}  
	
	
	\question P133页, 4.15题.
	设$E=\{\langle M\rangle|M\text{是一个DFA，它接受1比0多的串\}，证明}E$是可判定的。（提示 ：一些关于CFL的定理对此会有帮助。）
	\begin{solution}
		定义语言 
        \[
        C = \{ w \in \{0,1\}^* \mid \#_1(w) > \#_0(w) \}.
        \]
        已知 \(C\) 是上下文无关语言。  
        
        我们需判定：对于给定的 DFA \(M\)，是否 \(L(M) \subseteq C\)。
        
        等价条件：检查是否有 \(L(M) \cap \overline{C} = \varnothing\)，其中 \(\overline{C} = \{ w \mid \#_1(w) \le \#_0(w) \}\)。  
        \(\overline{C}\) 也是 CFL。
        
        构造判定器如下：
        
        \(D =\) “输入 \(\langle M \rangle\)，其中 \(M\) 是一个 DFA：
        \begin{enumerate}
            \item 构造一个 CFG \(G_1\) 使得 \(L(G_1) = \overline{C}\)。
            \item 由于 DFA 语言与 CFL 的交是 CFL，构造一个 CFG \(G_2\) 使得 \(L(G_2) = L(M) \cap L(G_1)\)。
            \item 在 \(G_2\) 上运行 \(E_{\text{CFG}}\) 的判定器，检查 \(L(G_2)\) 是否为空。
            \item 如果 \(E_{\text{CFG}}\) 判定器输出“空”，说明 \(L(M) \cap \overline{C} = \varnothing\)，即 \(L(M) \subseteq C\)，则接受 \(\langle M \rangle\)；否则拒绝。”
        \end{enumerate}
        \textbf{正确性分析}：
        \begin{itemize}
            \item 若 \(\langle M \rangle \in E\)，即 \(L(M) \subseteq C\)，则 \(L(M) \cap \overline{C} = \varnothing\)，所以 \(L(G_2) = \varnothing\)，\(E_{\text{CFG}}\) 判定器输出“空”，从而 \(D\) 接受。
            \item 若 \(\langle M \rangle \notin E\)，则存在 \(w \in L(M)\) 但 \(w \notin C\)，即 \(w \in \overline{C}\)，于是 \(w \in L(M) \cap \overline{C} = L(G_2)\)，因此 \(L(G_2) \neq \varnothing\)，\(E_{\text{CFG}}\) 判定器输出“非空”，从而 \(D\) 拒绝。
        \end{itemize}
        由于所有步骤（构造 CFG、交集 CFG、空性检查）都是可计算的，且 \(E_{\text{CFG}}\) 是可判定的，所以 \(D\) 是一个判定器。  
        因此，\(E\) 是可判定的。
	\end{solution}  
	
	\question P133页, 4.17题.
	设$BAL_\text{DFA}=\{\langle M\rangle|M$是一个DFA，它接受含有相同个数的0和1的串\}，证明$BAL_{DFA}$是可判定的。（提示 ：一些关于CFL的定理对此会有帮助。）
	\begin{solution}
		定义语言 
        \[
        B = \{ w \in \{0,1\}^* \mid \#_0(w) = \#_1(w) \}.
        \]
        已知 \(B\) 是上下文无关语言。
        
        题目等价为：
        \[
        BAL_{\text{DFA}} = \{ \langle M \rangle \mid M \text{ 是一个 DFA 并且 } L(M) \cap B \neq \varnothing \}.
        \]
        即 \(M\) 接受至少一个 \(0\) 和 \(1\) 个数相等的字符串。
        
        构造判定器如下：
        
        \(D =\) “输入 \(\langle M \rangle\)，其中 \(M\) 是一个 DFA：
        \begin{enumerate}
            \item 构造一个 CFG \(G_1\) 使得 \(L(G_1) = B\)。
            \item 由于 DFA 语言与 CFL 的交是 CFL，构造一个 CFG \(G_2\) 使得 \(L(G_2) = L(M) \cap L(G_1)\)。
            \item 在 \(G_2\) 上运行 \(E_{\text{CFG}}\) 的判定器，检查 \(L(G_2)\) 是否为空。
            \item 如果 \(E_{\text{CFG}}\) 判定器输出“非空”，则接受 \(\langle M \rangle\)（因为存在 \(w \in L(M) \cap B\)）；否则拒绝。”
        \end{enumerate}
        
        \textbf{正确性分析}：
        \begin{itemize}
            \item 若 \(\langle M \rangle \in BAL_{\text{DFA}}\)，则存在 \(w \in L(M) \cap B\)，从而 \(L(G_2) \neq \varnothing\)，\(E_{\text{CFG}}\) 判定器输出“非空”，于是 \(D\) 接受。
            \item 若 \(\langle M \rangle \notin BAL_{\text{DFA}}\)，则 \(L(M) \cap B = \varnothing\)，即 \(L(G_2) = \varnothing\)，\(E_{\text{CFG}}\) 判定器输出“空”，于是 \(D\) 拒绝。
        \end{itemize}
        
        以上所有步骤（构造 CFG、交集 CFG、空性检查）都是可计算的，且 \(E_{\text{CFG}}\) 是可判定的，因此 \(D\) 是一个判定器。  
        
        故 \(BAL_{\text{DFA}}\) 是可判定的。
	\end{solution}  
	
	
	
	\question P133页, 4.20题.
	设$S=\{\langle M\rangle|M\text{是一个DFA，且只要接受}w\text{，就接收}w^R\}$。证明 $S$ 是可判定的。
	\begin{solution}
		给定 DFA \(M = (Q, \Sigma, \delta, q_0, F)\)。  
        构造接受 \(L(M)\) 的逆序语言的 DFA \(M_R\) 方法如下：
        
        首先构造一个 NFA \(N\)：
        \[
        N = (Q \cup \{q_{\text{new}}\}, \Sigma, \delta_N, q_{\text{new}}, \{q_0\}),
        \]
        其中转移函数 \(\delta_N\) 定义如下：
        \begin{itemize}
            \item 对于原 DFA 的每个转移 \(\delta(p,a)=q\)，在 \(N\) 中添加 \(\delta_N(q,a)=p\)。
            \item 对每个接受状态 \(f \in F\)，添加 \(\epsilon\)-转移 \(q_{\text{new}} \to f\)（为了从新开始状态能到原来 DFA 的接受状态）。
        \end{itemize}
        这样 \(N\) 接受的语言恰好是 \(L(M)^R\)。  
        将 \(N\) 通过子集构造法确定化为 DFA \(M_R\)。
        
        于是，条件“只要接受 \(w\)，就接受 \(w^R\)”等价于要求 \(L(M) \subseteq L(M_R)\)。  
        但由于逆操作两次回到原串，若 \(w \in L(M_R)\)，则 \(w^R \in L(M)\)，交换 \(M\) 与 \(M_R\) 的角色可知 \(L(M_R) \subseteq L(M)\)，因此条件等价于 \(L(M) = L(M_R)\)。
        
        判断两个 DFA 是否等价是可判定的（可通过同步乘积构造，检查是否所有可达状态对的接受性一致）。  
        
        因此判定器 \(D\) 如下工作：
        \[
        D = \text{“输入 } \langle M \rangle:
        \]
        \begin{enumerate}
            \item 根据上述方法构造 DFA \(M_R\) 使得 \(L(M_R) = (L(M))^R\)。
            \item 判断 \(M\) 与 \(M_R\) 是否等价（使用 DFA 等价性判定算法）。
            \item 若等价，则接受；否则拒绝。”
        \end{enumerate}
        
        以上每一步都可在有限时间内完成，因此 \(S\) 是可判定的。
	\end{solution}  
	
	\question P134页, 4.23题。证明可判定的语言类在同态下不封闭.
	\begin{solution}
    
		通过构造反例来证明：存在一个可判定语言 $L$ 和一个同态 $h$，使得 $h(L)$ 不是可判定的。

        设 $A_{\text{TM}} = \{ \langle M, w \rangle \mid M \text{ 是图灵机且 } M \text{ 接受 } w \}$，这是著名的不可判定语言。
        
        考虑如下语言：
        \[
        L = \left\{ \langle M, w \rangle \# C \,\middle|\, 
        \begin{array}{l}
        M \text{ 是图灵机},\ w \in \{0,1\}^*, \\
        C \text{ 是 } M \text{ 在输入 } w \text{ 上的一个接受计算历史}
        \end{array}
        \right\}.
        \]
        
        这里，$\#$ 是一个不在 $\{0,1\}$ 中的新符号，而计算历史 $C$ 是对 $M$ 在 $w$ 上运行过程中从初始格局到接受格局的完整、合法的格局序列的编码（例如用特殊分隔符连接各格局）。只需检查：
        \begin{enumerate}
            \item 格局序列以 $M$ 在 $w$ 上的初始格局开始；
            \item 每个格局到下一个格局的转换符合 $M$ 的转移函数；
            \item 最后一个格局是接受格局。
        \end{enumerate}
        这些检查均可在有限步内完成，无需模拟 $M$ 的无限运行。因此，存在图灵机判定 $L$，即 $L$ 是可判定的。
        
        现在定义同态 $h: (\{0,1,\#\} \cup \Gamma)^* \to \{0,1\}^*$（其中 $\Gamma$ 是用于编码计算历史的字母表，与 $\{0,1\}$ 不交），满足：
        \[
        h(a) =
        \begin{cases}
        a & \text{若 } a \in \{0,1\}, \\
        \varepsilon & \text{若 } a \in \{\#\} \cup \Gamma.
        \end{cases}
        \]
        即 $h$ 保留输入中的 $\langle M, w \rangle$ 部分（由 0 和 1 构成），并删除 $\#$ 及其后的整个计算历史 $C$。
        
        于是，
        \[
        h(L) = \{ \langle M, w \rangle \mid \exists C \text{ 使得 } \langle M, w \rangle \# C \in L \}
              = \{ \langle M, w \rangle \mid M \text{ 接受 } w \}
              = A_{\text{TM}}.
        \]
        
        由于 $A_{\text{TM}}$ 是不可判定的，而 $L$ 是可判定的，这表明存在一个同态 $h$ 使得 $h(L)$ 不可判定。
        
        因此，可判定的语言类在同态下不封闭。
	\end{solution}
	
	
	
	
	\question P134页, 4.24.设$C$是一个语言。 证明$C$是图灵可识别的， 当且仅当存在一个可判定语言$D$，使得$C=\{x|\exists y(\langle x,y\rangle\in D)\}$
	\begin{solution}
    
        \noindent\textbf{($\rightarrow$) 若 $C$ 是图灵可识别的，则存在可判定语言 $D$ 使得 $C = \{x \mid \exists y\, (\langle x,y\rangle \in D)\}$.}
        
        因为 $C$ 是图灵可识别的，存在图灵机 $M$ 使得 $L(M) = C$。  
        构造语言
        \[
        D = \{ \langle x, t \rangle \mid M \text{ 在输入 } x \text{ 上于 } t \text{ 步内接受} \},
        \]
        其中 $t$ 以一元（如 $0^t$）或二进制形式编码为字符串 $y$，$\langle x, t \rangle$ 表示标准配对编码。
        
        我们断言 $D$ 是可判定的。给定输入 $\langle x, t \rangle$，图灵机 $N$ 可模拟 $M$ 在 $x$ 上运行至多 $t$ 步：
        - 若 $M$ 在 $t$ 步内进入接受状态，则 $N$ 接受；
        - 否则（包括拒绝或未停机），$N$ 拒绝。
        
        由于模拟步数有界（为 $t$），$N$ 总会停机，故 $D$ 是可判定的。
        
        此外，显然
        \[
        x \in C \iff M \text{ 接受 } x \iff \exists t\, (\langle x, t \rangle \in D),
        \]
        因此 $C = \{ x \mid \exists y\, (\langle x, y \rangle \in D) \}$。
        
        \noindent\textbf{($\leftarrow$) 若存在可判定语言 $D$ 使得 $C = \{x \mid \exists y\, (\langle x,y\rangle \in D)\}$，则 $C$ 是图灵可识别的.}
        
        设 $D$ 是可判定的，令 $M_D$ 为判定 $D$ 的图灵机。  
        构造图灵机 $M$ 识别 $C$ 如下：
        
        \begin{quote}
        \textbf{输入} $x$：\\
        枚举所有可能的字符串 $y_1, y_2, y_3, \dots$（按长度和字典序）；\\
        对每个 $i = 1, 2, 3, \dots$，在第 $i$ 阶段：\\
        \quad 模拟 $M_D$ 在 $\langle x, y_1 \rangle, \langle x, y_2 \rangle, \dots, \langle x, y_i \rangle$ 上运行各至多 $i$ 步；\\
        \quad 若任一模拟接受，则 $M$ 接受 $x$。
        \end{quote}
        
        更简单地说，$M$ 可以逐一枚举所有 $y$，并在每枚举一个 $y$ 后运行 $M_D$ 判定 $\langle x, y \rangle \in D$ 是否成立。一旦发现某个 $y$ 使得 $\langle x, y \rangle \in D$，就接受 $x$。
        
        由于 $D$ 是可判定的，对每个 $y$，$M_D$ 总会停机并给出正确答案。因此，若存在某个 $y$ 使得 $\langle x, y \rangle \in D$，则 $M$ 最终会找到它并接受；若不存在这样的 $y$，则 $M$ 永不停机。符合图灵可识别定义。
        
        故 $C$ 是图灵可识别的。
        
        \medskip
        
        综上，$C$ 是图灵可识别的当且仅当存在可判定语言 $D$ 使得 $C = \{ x \mid \exists y\, (\langle x, y \rangle \in D) \}$。
	\end{solution}
	
	\question P134页, 4.26题.
	设$A=\{\langle R
	\rangle|R \text{是一个正则表达式，其所描述的语言中至少有一个串}w\text{以}111$ $\text{为子串（即有 } x\text{ 和 }y\text{ 使得 }w=x111y\text{）}\}$。证明$A$是可判定的。
	
	\begin{solution}
        通过构造一个判定 $A$ 的图灵机来证明其可判定性。

        首先，注意到“包含子串 $111$”的语言是正则的。定义正则语言
        \[
        L_{111} = \{ w \in \{0,1\}^* \mid w \text{ 包含 } 111 \text{ 作为子串} \} = \Sigma^* 111 \Sigma^*,
        \]
        其中 $\Sigma = \{0,1\}$。显然，$L_{111}$ 可由正则表达式 $\Sigma^*111\Sigma^*$ 描述，因此是正则语言。

        \[
        \langle R \rangle \in A \iff \exists w \in L(R) \text{ 使得 } w \text{ 包含 } 111 \text{ 作为子串}
        \iff L(R) \cap L_{111} \neq \varnothing.
        \]
        
        因此给定输入 $\langle R \rangle$，其中 $R$ 是一个正则表达式，判断 $L(R) \cap L_{111} \neq \varnothing$ 是否成立。  
        
        
        现在利用正则语言的如下性质：
        \begin{enumerate}
            \item 给定正则表达式 $R$，可有效构造等价的 DFA $M_R$（通过标准算法：RE $\to$ NFA $\to$ DFA）；
            \item 给定 DFA $M_1$ 和 $M_2$，可构造识别 $L(M_1) \cap L(M_2)$ 的 DFA $M_\cap$（通过乘积构造）；
            \item 给定 DFA $M$，可判定 $L(M) = \varnothing$ 是否成立（通过检查从起始状态是否可达任意接受状态）。
        \end{enumerate}
        
        基于以上，构造图灵机 $T$ 判定 $A$ 如下：
        
        \begin{quote}
        \textbf{输入}：$\langle R \rangle$，其中 $R$ 是正则表达式。
        
        \begin{enumerate}
            \item 构造识别 $L_{111} = \Sigma^*111\Sigma^*$ 的 DFA $M_{111}$（例如，一个 4 状态 DFA，依次读取三个连续的 1 后进入吸收接受状态）。
            \item 将 $R$ 转换为等价的 NFA $N_R$，再转换为等价的 DFA $M_R$。
            \item 构造 DFA $M$ 使得 $L(M) = L(M_R) \cap L(M_{111})$（使用乘积构造）。
            \item 检查 $L(M)$ 是否为空：若为空，则拒绝；否则，接受。
        \end{enumerate}
        \end{quote}
        
        由于上述每一步都是可计算的且总在有限步内停机，图灵机 $T$ 总会停机并给出正确答案。因此，$A$ 是可判定的。
	\end{solution}
	
	
	\question P134页, 4.27题.
	证明确定一个 CFG 是否生成$1\star$ 中的所有串的问题是可判定的 。换言之 ， 证明$\{\langle G
	\rangle |G\text{是}\{0,1\}\text{ 上的一个CFG，且 }1\star\subseteq L(G)\}$ 是一个可判定语言。
	\begin{solution}
		检查 $L(G')$ 是否包含所有长度不超过某个界 $N$ 的 $1^n$。

        具体地，设 $G'$ 有 $k$ 个变量。在 CNF 下，若 $1^n \in L(G')$ 且 $n > 2^k$，则 $1^n$ 可泵。  
        因此，若 $G'$ 缺少某个 $1^n$，则必存在一个最小的缺失长度 $n_0 \leq 2^k$（否则，若所有 $1^m$ 对 $m \leq 2^k$ 都在 $L(G')$ 中，则通过泵引理可生成所有更长的 $1^n$）。
        
        于是，我们得到一个有限检验算法：
        
        \begin{quote}
        \textbf{判定算法：}
        
        输入：CFG $G$（字母表含 $\{0,1\}$）。
        
        \begin{enumerate}
            \item 构造 CFG $G'$ 使得 $L(G') = L(G) \cap 1^*$（通过与 DFA for $1^*$ 交）。
            \item 将 $G'$ 转换为 Chomsky 范式，设其变量数为 $k$，令 $N = 2^k$。
            \item 对每个 $n = 0, 1, 2, \dots, N$，判定 $1^n \in L(G')$ 是否成立（CFG 的成员问题是可判定的，可用 CYK 算法）。
            \item 若所有 $1^n$（$0 \leq n \leq N$）都属于 $L(G')$，则接受；否则拒绝。
        \end{enumerate}
        \end{quote}
	\end{solution}  
	
	
	\question P134页, 4.29题.
	设$A=\{\langle R,S\rangle |R\text{和}S\text{是正则表达式，且}L(R)\subseteq L(S)\}$，证明$A$是可判定的。
	\begin{solution}
		利用正则语言的封闭性质和若干已知的可判定问题来构造一个判定 $A$ 的图灵机。

        集合包含关系的等价转换：
        \[
        L(R) \subseteq L(S) \quad \Longleftrightarrow \quad L(R) \cap \overline{L(S)} = \varnothing.
        \]
        
        由于正则语言在补运算和交运算下封闭，且空性问题是可判定的，我们可以据此设计判定算法。
        
        具体步骤如下：
        
        \begin{quote}
        \textbf{输入}：$\langle R, S \rangle$，其中 $R$ 和 $S$ 是正则表达式。
        
        \begin{enumerate}
            \item 将正则表达式 $R$ 和 $S$ 分别转换为等价的 DFA $M_R$ 和 $M_S$。  
                  （这是可行的：通过标准算法 $ \text{RE} \to \text{NFA} \to \text{DFA} $。）
        
            \item 构造 DFA $M_{\overline{S}}$ 使得 $L(M_{\overline{S}}) = \overline{L(S)}$。  
                  （只需将 $M_S$ 的接受状态与非接受状态互换。）
        
            \item 构造 DFA $M$ 使得 $L(M) = L(M_R) \cap L(M_{\overline{S}})$。  
                  （使用乘积构造法：状态为 $(q_R, q_{\overline{S}})$，接受状态为两者均为接受的状态。）
        
            \item 判定 $L(M)$ 是否为空：  
                  从起始状态出发，执行图搜索（如 BFS 或 DFS），检查是否存在一条路径到达某个接受状态。  
                  若存在，则 $L(M) \neq \varnothing$；否则 $L(M) = \varnothing$。
        
            \item 若 $L(M) = \varnothing$，则接受（即 $L(R) \subseteq L(S)$）；否则拒绝。
        \end{enumerate}
        \end{quote}
        
        上述每一步都是有效可计算的，并且总在有限步内停机。
        因此，存在一台图灵机总能正确判定 $\langle R, S \rangle \in A$ 是否成立，故 $A$ 是可判定的。
		
	\end{solution}  
	
	\question P134页, 4.31题.
	设$INFINITE_\text{PDA}=\{\langle M\rangle|M\text{是一个PDA，且}L(M)\text{是一个无限语言}\}$。证明$INFINITE_\text{PDA}$是可判定的。
	\begin{solution}
		通过将 PDA 转换为等价的上下文无关文法（CFG），并利用上下文无关语言的泵引理来设计一个判定算法。
        
        \noindent\textbf{判定算法如下：}
        
        \begin{quote}
        \textbf{输入}：$\langle M \rangle$，其中 $M$ 是一个 PDA。
        
        \begin{enumerate}
            \item 将 PDA $M$ 转换为等价的上下文无关文法 $G$，使得 $L(G) = L(M)$。  
                  （这是可行的：Sipser 定理 2.22 给出了从 PDA 到 CFG 的有效构造。）
        
            \item 将 $G$ 转换为 Chomsky 范式（CNF）。设 $G$ 中有 $k$ 个变量。  
                  （转换到 CNF 是机械过程，且不改变语言。）
        
            \item 令泵长度 $p = 2^k$。（在 CNF 下，任何派生树深度超过 $k$ 必有重复变量，故 $p = 2^k$ 是有效泵长度。）
        
            \item 对每个长度 $n$ 满足 $p \leq n < 2p$，枚举所有长度为 $n$ 的字符串 $w \in \Sigma^n$，并使用 CYK 算法（或其他 CFG 成员性判定算法）检查是否 $w \in L(G)$。
        
            \item 若存在某个 $w$（长度在 $[p, 2p)$ 内）使得 $w \in L(G)$，则接受（即 $L(M)$ 无限）；  
                  否则，拒绝（即 $L(M)$ 有限）。
        \end{enumerate}
        \end{quote}
        
        \noindent\textbf{正确性说明：}
        
        - 若 $L(M)$ 是无限的，则根据泵引理，必存在某个 $w \in L(M)$ 满足 $|w| \geq p$。  
          考虑最短的这样的 $w$。若 $|w| \geq 2p$，则其派生树中必有可泵子树，从而可“缩短”得到一个更短但仍 $\geq p$ 的字符串仍在 $L(G)$ 中，矛盾于最短性。  
          因此，最短的 $w$ 满足 $p \leq |w| < 2p$，会被算法找到。
        
        - 反之，若算法找到某个 $w$ 满足 $p \leq |w| < 2p$ 且 $w \in L(G)$，则由泵引理，$w$ 可被泵出无穷多个不同长度的字符串，故 $L(G)$（即 $L(M)$）是无限的。
        
        - 若对所有 $n \in [p, 2p)$ 都没有 $w \in L(G)$，则 $L(G)$ 中所有字符串长度均 $< p$，故 $L(G)$ 是有限集。

        整个过程在有限步内完成，并总是停机。因此，存在图灵机判定 $\mathrm{INFINITE}_{\text{PDA}}$，故该语言是可判定的。
		
	\end{solution}  
	
\end{questions}

下面是选做题。

\begin{questions}
	
	\question P134页, 4.28题
	设$\Sigma=\{0,1\}$，证明确定一个CFG是否生成$1\star$中的串的问题是可判定的。换言之，证明$\{\langle G\rangle|G\text{是}\{0,1\}\text{上的一个CFG，且}1\star \cap L(G)\neq \emptyset\}$是一个可判定语言。
	\begin{solution}
		
		
	\end{solution}  
	
	\question P134页, 4.30题
	设$A=\{\langle M \rangle|M\text{是DFA，它不接受任何包含奇数个1的串}\}$，证明$A$是可判定的。
	\begin{solution}
		
		
	\end{solution}   
	
	
	\question P134页, 4.32题
	设$INFINITE_\text{DFA}=\{\langle A\rangle|A\text{是一个DFA，且}L(A)\text{是一个无限语言}\}$。证明$INFINITE_\text{DFA}$是可判定的。
	\begin{solution}
		
		
	\end{solution}  
	
	
\end{questions}
 	
\end{document}
